% !TeX encoding = UTF-8
\documentclass[variant=guia,cover=institutional,mode=screen]{ua-engineering-v6-article}
% !TeX encoding = UTF-8
\UASetUniversity{Universidad Austral}
\UASetFaculty{Facultad de Ingeniería}
\UASetDepartment{Departamento de Computación}
\UASetCampus{Pilar}
\UASetCareer{Ingeniería Informática}
\UASetCommission{Comisión A}
\UASetProfessor{Equipo docente}
\UASetSemester{Segundo cuatrimestre 2026}
\UASetCourse{Sistemas Distribuidos}
\UASetDate{2026}


\UASetClassNumber{15}
\UASetClassTitle{Diseño integral de sistemas distribuidos}
\UASetDocKind{Guía de ejercicios}

\title{Clase 15 --- Guía de ejercicios}
\author{\UAProfessor}
\date{\UADate}

\begin{document}

\section{Indicaciones generales}
Resuelva cada ejercicio declarando supuestos, modelo de fallas, garantía y trade-off. Cuando corresponda, construya una ejecución verificable y vincule la respuesta con evidencia observable.

\section{Nivel 1: fundamentos y reconocimiento}
\begin{uaexercise}
Explique por qué diseñar un sistema distribuido no equivale a listar tecnologías o patrones de moda.
\end{uaexercise}

\begin{uaexercise}
Explique el sentido del esquema $D=(P,F,G,S,T)$ y por qué ayuda a estructurar una defensa arquitectónica.
\end{uaexercise}

\begin{uaexercise}
Explique por qué una buena propuesta debe hacer explícitas las fallas asumidas y no solo el camino feliz.
\end{uaexercise}

\begin{uaexercise}
Explique por qué una arquitectura técnicamente seria debe reconocer también lo que no resuelve.
\end{uaexercise}

\begin{uaexercise}
Compare consistencia fuerte y consistencia eventual como decisiones de diseño, no como slogans.
\end{uaexercise}


\section{Nivel 2: ejecuciones y fallas}
\begin{uaexercise}
Compare monolito modular y microservicios en términos de complejidad total del sistema.
\end{uaexercise}

\begin{uaexercise}
Compare integración por RPC e integración por eventos en términos de acoplamiento, trazabilidad y operación.
\end{uaexercise}

\begin{uaexercise}
Compare transacción fuerte y saga en términos de garantías, costo y semántica visible.
\end{uaexercise}

\begin{uaexercise}
Diseñe una estrategia de defensa para responder la objeción: “¿por qué no un monolito?”.
\end{uaexercise}

\begin{uaexercise}
Diseñe una estrategia de defensa para responder la objeción: “¿por qué no consistencia fuerte en todo?”.
\end{uaexercise}


\section{Nivel 3: diseño y comparación}
\begin{uaexercise}
Diseñe una estrategia de defensa para responder la objeción: “¿qué pasa bajo partición?”.
\end{uaexercise}

\begin{uaexercise}
Diseñe una estrategia de defensa para responder la objeción: “¿cómo observás y operás este sistema?”.
\end{uaexercise}

\begin{uaexercise}
Diseñe una arquitectura razonable para un e-commerce global y justifique sus decisiones principales.
\end{uaexercise}

\begin{uaexercise}
Diseñe una arquitectura razonable para una plataforma financiera global y justifique sus decisiones principales.
\end{uaexercise}

\begin{uaexercise}
Diseñe una arquitectura razonable para una plataforma de streaming de eventos y justifique sus decisiones principales.
\end{uaexercise}


\section{Nivel 4: integración y defensa}
\begin{uaexercise}
Compare por qué esos tres casos no deberían resolverse con exactamente la misma arquitectura.
\end{uaexercise}

\begin{uaexercise}
Explique cómo deberían aparecer observabilidad y SLOs dentro de una defensa arquitectónica seria.
\end{uaexercise}

\begin{uaexercise}
Explique cómo deberían aparecer resiliencia e incident management dentro de una defensa seria.
\end{uaexercise}

\begin{uaexercise}
Explique por qué ownership de datos y bounded contexts no son detalles secundarios.
\end{uaexercise}

\begin{uaexercise}
Explique por qué una propuesta distribuida sin modelo operacional está incompleta.
\end{uaexercise}


\end{document}
